\documentclass[a4paper,11pt]{article}

\usepackage[
    a4paper,
    top=2.54cm,
    bottom=2.54cm,
    left=2.54cm,
    right=2.54cm
]{geometry}

\usepackage{xcolor}
\usepackage{tcolorbox}
\usepackage{enumitem}
\usepackage{hyperref}
\usepackage{titlesec}
\usepackage{booktabs}
\usepackage{tabularx}
\usepackage{array}

\setlist[itemize]{noitemsep, topsep=2pt}
\setlist[enumerate]{noitemsep, topsep=2pt}

\titleformat{\section}
    {\large\bfseries}{\thesection.}{0.5em}{}

\titleformat{\subsection}
    {\normalsize\bfseries}{\thesubsection.}{0.5em}{}

% Placeholder for values to be filled in by students.
\newcommand{\blank}[1][1.2cm]{\rule{#1}{0.4pt}}

\begin{document}

\begin{tcolorbox}[
    colback=white,
    colframe=black,
    width=\textwidth,
    arc=3mm,
    boxrule=1pt
]
\begin{center}
    \Large \textbf{Efficient Computing for Artificial Intelligence}\\[3pt]
    \large \textbf{Homework 1}\\[5pt]
    \textbf{\colorbox{yellow!30}{DUE DATE: 2 Dec (h23:59)}}\\[8pt]
\end{center}

\textbf{Submission Instructions:}\\[3pt]

Each group must send an e-mail to
\texttt{daniele.jahier@polito.it}, with
\texttt{valentino.peluso@polito.it} in CC.
Use the subject
\textbf{ECAI27 Group\texttt{N}}, replacing \texttt{N} with
the group ID.\\[3pt]

Attach a single ZIP archive named
\textbf{HW1\_Group\texttt{N}.zip}, containing:
\begin{enumerate}
    \item The code deliverables specified in each exercise.
    \item The PDF report.
\end{enumerate}

Late submissions or submissions that do not comply with these
instructions will be \textbf{automatically discarded}.
\end{tcolorbox}

\twocolumn[
\begin{@twocolumnfalse}
\begin{center}
\Large \textbf{HW1 – Group\texttt{N}}\\
\normalsize (Replace \texttt{N} with your group ID)\\

\begin{tabular}{lll}
\textbf{Surname Name} (s\texttt{XXXXXX}) & \textbf{Surname Name} (s\texttt{XXXXXX}) & \textbf{Surname Name} (s\texttt{XXXXXX}) \\
\end{tabular}
\end{center}
\end{@twocolumnfalse}
]

\section*{Results}

Complete the table with the operational intensity (OI), median latency,
and standard deviation of latency for each kernel.

\begin{itemize}
    \item \textbf{Operational intensity:}
    Report the ideal OI in FLOP/byte, assuming a single-level memory
    hierarchy and an infinite number of registers. 
    \item \textbf{Latency:}
    Report the median and standard deviation of the samples in milliseconds (ms), rounded
    to one decimal digit.
\end{itemize}

\begin{center}
\renewcommand{\arraystretch}{1.25}
\begin{tabularx}{\textwidth}{
    @{}
    >{\raggedright\arraybackslash}X
    >{\centering\arraybackslash}p{2.0cm}
    >{\centering\arraybackslash}p{2.6cm}
    >{\centering\arraybackslash}p{2.6cm}
    @{}
}
\toprule
Kernel
& OI (FLOP/byte)
& Median latency (ms)
& Std. latency (ms) \\
\midrule
Dense scalar
& \blank
& \blank[1.5cm]
& \blank[1.5cm] \\

Dense tiled with NEON and OpenMP
& \blank
& \blank[1.5cm]
& \blank[1.5cm] \\

Sparse scalar
& \blank
& \blank[1.5cm]
& \blank[1.5cm] \\

Sparse NEON
& \blank
& \blank[1.5cm]
& \blank[1.5cm] \\

Sparse tiled
& \blank
& \blank[1.5cm]
& \blank[1.5cm] \\

Sparse OpenMP (2 threads)
& \blank
& \blank[1.5cm]
& \blank[1.5cm] \\

Sparse OpenMP (4 threads)
& \blank
& \blank[1.5cm]
& \blank[1.5cm] \\

Sparse Combined (2 threads)
& \blank
& \blank[1.5cm]
& \blank[1.5cm] \\

Sparse Combined (4 threads)
& \blank
& \blank[1.5cm]
& \blank[1.5cm] \\

\bottomrule
\end{tabularx}
\end{center}

\end{document}